\section{Introduction}\label{sec:introduction}

Modern machine learning systems depend on data infrastructure that can store
features, update them from new events, and recover from failures. The goal of
MLStore-Lite is to build a small educational prototype of such a system. The
project is inspired by selected parts of \textit{Designing Data-Intensive
Applications} (DDIA), but keeps the implementation local and readable.

In concrete terms, the final project is a local ML infrastructure prototype. It
contains a storage engine, leader-follower replication, sharding, batch
processing, stream processing, observability, online inference, a sequential
recommender, metadata, lineage, and data quality checks.

The project does not try to compete with production tools such as RocksDB,
Cassandra, Kafka, Spark, Flink, or managed ML platforms. I use those systems as
reference points, not as tools I claim to have mastered deeply. The purpose is
to understand their architecture by implementing compact versions of central
ideas: write-ahead logging, SSTable-like files, replication, sharding, batch
processing, stream processing, observability, feature serving, model inference,
sequential recommendation, metadata, lineage, and data quality.

The project simulates distributed ideas locally. It does not implement real
network communication, consensus, automatic leader election, or production-grade
fault tolerance. The AI layer is included not as a standalone model project, but
to show why data infrastructure matters for ML systems: models need features,
serving paths, logs, reproducibility, and traceability.

This project is also personal learning work. Coming from a mathematics
background and having worked with code for a bit less than two years, I wanted
to go deeper into software architecture rather than only use high-level ML
libraries. The implementation was AI-assisted using Codex 5.5, Claude Sonnet
4.7, Claude Opus 4.7, and experiments with local LLMs such as Qwen models around
35B parameters. The tools helped with implementation, debugging, and writing,
but the project scope and interpretation remained part of the learning process.

\section{Reading Scope and Method}\label{sec:method}

The reading scope focused on selected DDIA chapters rather than the whole book.
The main mapping was:

\begin{center}
\footnotesize
\setlength{\tabcolsep}{4pt}
\renewcommand{\arraystretch}{0.95}
\begin{tabular}{p{0.22\linewidth} p{0.32\linewidth} p{0.36\linewidth}}
\toprule
\textbf{Milestone} & \textbf{DDIA connection} & \textbf{Implementation} \\
\midrule
Storage & Ch. 3: Storage and Retrieval & WAL, memtable, SSTable-like files, compaction, recovery \\
Replication & Ch. 5: Replication & leader-follower replication, sync/async writes, failover \\
Sharding & Ch. 6: Partitioning & consistent hashing, virtual nodes, routing, rebalancing \\
Batch & Ch. 10: Batch Processing & local map-shuffle-reduce engine and batch features \\
Stream & Ch. 11: Stream Processing & event log, producers, consumers, offsets, windows \\
Integration & system composition & one object wiring storage, replication, sharding, batch, stream \\
Observability & operational tradeoffs & structured logs, timings, JSON-lines experiment records \\
AI extension & ML infrastructure & feature serving, inference, sequential recommender \\
Inspection layer & MLOps traceability & feature registry, quality checks, lineage records \\
\bottomrule
\end{tabular}
\vspace{0.2cm}
\textit{Main reading-to-implementation mapping.}
\end{center}

Several topics were intentionally not implemented deeply: transactions,
consensus, automatic leader election, distributed commit, real network
partitions, and production recovery. This boundary keeps the project focused on
understanding the architecture instead of reproducing a full database or stream
processor.

The AI/recommender part is designed around ecommerce event data. The code can
read the RetailRocket ecommerce dataset from Kaggle if the file is placed at
\path{data/raw/retailrocket/events.csv}. That dataset has event-style columns
such as timestamp, visitor id, event type, and item id \cite{retailrocketdataset}.
The repository does not commit the Kaggle CSV because it is external and larger
than the source code. When the file is missing, the scripts use a small built-in
sample with the same shape: user id, event type, item id, and timestamp. This
keeps the project reproducible while still showing where a larger real dataset
fits into the architecture.

\section{System Architecture}\label{sec:architecture}

MLStore-Lite is organized as a layered system:

\begin{lstlisting}
storage -> replication -> sharding -> batch -> stream
        -> integration -> observability -> inference
        -> sequential recommender -> metadata/lineage/quality
\end{lstlisting}

The storage layer persists key-value data on one node. The replication layer
keeps several copies of the same data. The sharding layer decides which
replicated group owns each key. Batch and stream processing turn raw events into
features. The AI layers consume those features and event histories. The final
inspection layer explains feature meanings, validates input events, and records
which data was used for predictions.

\begin{figure}[H]
\centering
\begin{tikzpicture}[
    font=\scriptsize,
    node distance=0.65cm and 0.55cm,
    >={Stealth[length=2mm]},
    stage/.style={rounded corners=4pt, thick, align=center, minimum width=2.05cm, minimum height=0.72cm},
    blue/.style={stage, draw=flowBlue!85!black, fill=flowBlue!10},
    green/.style={stage, draw=flowGreen!85!black, fill=flowGreen!10},
    gold/.style={stage, draw=flowGold!85!black, fill=flowGold!12},
    rose/.style={stage, draw=flowRose!85!black, fill=flowRose!10},
    ink/.style={stage, draw=flowInk!85!black, fill=flowInk!8},
    arrow/.style={->, thick, draw=flowInk!75},
    splitarrow/.style={->, thick, draw=flowBlue!75},
    note/.style={align=center, text=flowInk!80}
]
    \node[blue] (events) at (0,0) {Events\\Kaggle or sample};
    \node[green, right=of events] (quality) {Quality\\checks};
    \node[gold, above right=0.35cm and 0.7cm of quality] (batch) {Batch\\features};
    \node[gold, below right=0.35cm and 0.7cm of quality] (stream) {Stream\\windows};
    \node[blue, right=2.45cm of quality] (store) {Sharded replicated\\feature store};
    \node[rose, above right=0.35cm and 0.75cm of store] (inference) {Online\\inference};
    \node[rose, below right=0.35cm and 0.75cm of store] (recommender) {Sequential\\recommender};
    \node[ink, right=2.55cm of store] (audit) {Prediction logs\\lineage + metrics};

    \draw[arrow] (events) -- (quality);
    \draw[splitarrow] (quality) -- (batch);
    \draw[splitarrow] (quality) -- (stream);
    \draw[arrow] (batch) -- (store);
    \draw[arrow] (stream) -- (store);
    \draw[arrow] (store) -- (inference);
    \draw[arrow] (store) -- (recommender);
    \draw[arrow] (inference) -- (audit);
    \draw[arrow] (recommender) -- (audit);

    \node[note, below=1.15cm of store] (core) {DDIA-style internals: WAL, SSTable-like files, replication, sharding, logs, offsets};
    \draw[arrow, dashed] (core) -- (store);
\end{tikzpicture}
\caption{Final architecture of MLStore-Lite. Events become validated inputs, features, model predictions, and traceable logs.}
\end{figure}

\subsection{Data Model and Key Design}

The common data model is deliberately simple: the system stores keys and values.
That choice makes the lower layers generic. The storage engine does not need to
know whether a value is a click count, a model probability, or a lineage record.
Meaning is added by the layers above it through naming conventions and metadata.

Most feature keys follow a readable pattern:

\begin{lstlisting}
feature:user:<user_id>:<feature_name>
\end{lstlisting}

For example, \texttt{feature:user:42:click\_count} means that the value belongs
to user 42 and stores a derived click-count feature. Prediction keys and lineage
records use the same idea: the prefix tells the reader what kind of object is
stored, while the rest of the key identifies the user, model, or run. This is
not only a coding convenience. In a partitioned system, the key is also the
input to the shard router. Therefore, key design affects both human
understanding and physical data placement.

\subsection{How the Local System Is Set Up}

The project simulates a small cluster on one machine. A node is a local storage
engine backed by a directory. A replica is another node directory containing a
copy of the same shard data. A shard is a logical partition of the key space,
not a single file. The default topology used in the demos is:

\begin{lstlisting}
2 shards
3 replicas per shard
6 local node directories
\end{lstlisting}

This local topology keeps the project runnable on a laptop while preserving the
important architecture: keys are routed to shards, each shard has multiple
replicas, and higher layers do not need to know which physical directory stores
the final value.

\section{Building the System Step by Step}\label{sec:milestones}

The project was implemented week by week. Each week added one system idea while
still reusing the layers before it. This was useful because it made the
architecture easier to reason about: a later layer never appears as an isolated
script, but as another piece on top of the storage and data movement layers.

\subsection{Week 1: Storage Engine}

The storage engine uses a write-ahead log (WAL), memtable, immutable SSTable-like
files, and compaction. A write is appended to the WAL before in-memory state is
updated, so the node can recover after restart. Recent writes live in the
memtable, and older flushed data lives in sorted files. Compaction merges those
files and removes overwritten values.

The DDIA idea behind this part is that a database is not only a dictionary from
keys to values. It also needs a physical strategy for writing to disk, recovering
after crashes, and keeping lookup cost reasonable as data grows. The WAL gives
durability, the memtable gives fast recent writes, and SSTable-like files make
disk data immutable and easier to merge. This is a simplified LSM-tree style
design, inspired by the same family of ideas used by systems such as LevelDB and
RocksDB.

\subsection{Week 2: Replication}

Replication uses a leader-follower model. Each replicated group has one leader
and two followers, giving replication factor three. Writes go to the leader and
are copied to followers. The project also supports manual failover by promoting
a follower.

This implementation does not claim to solve consensus. It is closer to a small
laboratory version of leader-based replication from DDIA: one node accepts the
write first, followers copy the result, and reads can still work after one node
is removed from service. The useful lesson is the separation between a logical
record and its physical copies. The value for a key is one logical fact, but the
system may store that fact on several nodes for fault tolerance.

\subsection{Week 3: Sharding}

Sharding divides different keys across replicated groups. MLStore-Lite uses
consistent hashing, where each key is hashed onto a ring and routed to the next
shard. This avoids manually assigning key ranges and makes rebalancing easier
when shards are added. In the demos there are two shards, and each shard has
three local replicas.

This is where partitioning and replication meet. A shard is not one file; it is
the ownership of a subset of keys. A replica is not a different partition; it is
another copy of the same shard data. In this project, a key such as
\texttt{feature:user:42:click\_count} is routed to one shard, and that shard is
stored by three local node directories. The consistent hash ring decides the
owner, while the replication layer decides how many copies exist inside that
owner group.

\subsection{Week 4: Batch Processing}

The batch engine follows a small MapReduce pattern:

\begin{lstlisting}
map -> shuffle -> reduce
\end{lstlisting}

Raw historical events are mapped into intermediate feature pairs, grouped by
feature key, and reduced into final values. Example stored keys are:

\begin{lstlisting}
feature:user:42:click_count
feature:user:42:purchase_count
\end{lstlisting}

The important idea is derived data. Raw events are the original observations,
while features are computed summaries that models can use. For example, many
old click and purchase events can become a small number of per-user counts. The
batch layer is useful when the input is finite: a file, a historical dataset, or
all events up to a chosen time. It is intentionally small, but it keeps the
central MapReduce shape from DDIA: map local records, shuffle by key, then reduce
each group.

\subsection{Week 5: Stream Processing}

The stream layer handles newer events. Producers append events to a local event
log. Consumers read records and store offsets. A stream processor groups events
into tumbling windows and writes recent feature updates to the same sharded
replicated store. This is a small local version of ideas found in Kafka, Flink,
and Kafka Streams.

The difference from batch processing is the role of time. In batch processing,
the dataset is already there. In stream processing, new records keep arriving,
and the system must remember how far it has read. This is why offsets matter:
an offset is the consumer's position in the event log. Tumbling windows then
group recent records into fixed time intervals before writing feature updates.

\subsection{Week 6--7: Integrated Feature Store}

After storage, replication, sharding, batch, and stream processing existed
separately, the integration layer connected them into one object. This matters
because infrastructure value comes from composition. A batch job should not only
print feature values; it should write them to the same store that stream
processing and inference will read from. The integrated feature store therefore
acts as the bridge between data engineering and ML serving.

The integrated layer also makes the read/write path more realistic. A client can
ask the system to write a feature without manually choosing a node directory.
The request goes through routing, lands on the correct replicated group, and is
then persisted by the underlying storage engine. That path is small in code, but
it contains the core idea behind many data platforms: separate the user-facing
API from the physical placement of data.

\subsection{Week 8: Evaluation and Observability}

Week 8 adds observability: structured logs, timing measurements, and JSON-lines
experiment records. The evaluation script measures batch feature runtime, stream
production, stream processing, feature reads, and shard distribution. These
measurements are not production benchmarks. They are local evidence that the
system runs and can be inspected.

This part connects the project to MLOps debugging and logging ideas. The goal is
not to create a monitoring platform. The goal is to make runs reproducible: if a
demo produces a timing number or a prediction, there is a small JSON-lines
record showing what happened, when it happened, and which parameters were used.
That makes the system easier to debug and easier to discuss in the report.

\subsection{Week 9: Online Inference}

Week 9 adds online feature serving and inference. A feature server reads stored
features for a user, a deterministic purchase-intent model computes a purchase
probability, and predictions are written to a JSON-lines log. The model also
reports confidence and warnings when expected features are missing.

The key architectural point is that the model does not own the data pipeline.
The model consumes features that previous layers have prepared. This makes the
AI layer a client of the infrastructure, not a replacement for it. In practical
ML systems this distinction is important: model quality depends not only on the
formula used for scoring, but also on whether the features are available,
consistent, and recent enough for the decision being made.

\subsection{Week 10: Scaling Experiments and Cloud Design}

Week 10 added local scaling experiments and a cloud architecture design note.
The workload experiment increases the number of events and records how runtime
changes. The shard hotspot experiment shows what happens when many keys route
to the same shard. These are still laptop-scale observations, but they make the
limitations visible instead of hiding them.

The cloud design document explains how the same conceptual layers could map to
managed services: durable object storage, managed streaming, distributed compute,
containerized services, observability, and model serving. Nothing in the project
requires cloud infrastructure, but the design note helps connect the local code
to realistic production architecture.

\subsection{Week 11: Sequential Recommendation}

Week 11 adds a small Transformer-style sequential recommender. Instead of only
using counts, it looks at ordered user behavior such as:

\begin{lstlisting}
view laptop -> view laptop -> add_to_cart laptop
\end{lstlisting}

Events and items become tokens, tokens become numeric IDs, a small attention
mechanism weighs recent history, and the model outputs a purchase probability.
This is not a production deep-learning system. It is an educational version of
sequence-based recommendation.

The intended larger input for this layer is the RetailRocket Kaggle ecommerce
events dataset. Its event rows can be converted into ordered user histories:
\texttt{visitorid} becomes the user, \texttt{event} becomes actions such as
view, add-to-cart, or purchase, \texttt{itemid} becomes the item token, and
\texttt{timestamp} orders the sequence. In the submitted repo, if that Kaggle
file is not present, the same code path uses the built-in sample so that the
training and demo commands still run on any machine.

\subsection{Week 12: Metadata, Lineage, and Data Quality}

Week 12 adds metadata, lineage, and data quality. The feature registry explains
stored feature keys. The quality layer validates events before processing. The
lineage log records which feature keys or event histories were used for a
prediction. The idea is simple:

\begin{lstlisting}
quality checks protect the input
feature registry explains the stored data
lineage explains the prediction
\end{lstlisting}

This final layer is not about making the model more complex. It is about making
the system more understandable. A prediction is easier to trust when the feature
definitions are documented, invalid input events are rejected, and the system can
show which features or event sequences contributed to the output.

\section{End-to-End Flow}\label{sec:flow}

The final system can be read as one pipeline. Historical events are processed by
the batch layer. Recent events are appended to the stream log and processed in
windows. Both paths write derived features into the same sharded, replicated
feature store. The inference layer reads those features, and the sequential
recommender can also read ordered event histories. Outputs are written as
prediction records and lineage records.

\begin{lstlisting}
historical events -> batch features
recent events     -> stream windows
features          -> sharded replicated store
store + histories -> inference and recommender
predictions       -> logs, lineage, quality reports
\end{lstlisting}

The important point is that the AI layer is downstream of the data system. A
model score is not just a mathematical function. It depends on whether events
were accepted by quality checks, whether features were computed, whether the
right shard served the data, and whether the prediction was logged for later
inspection. This is why the project combines DDIA-style infrastructure with
MLOps-style observability.

\section{Results and Reproducibility}\label{sec:results}

The project can be checked locally with one command:

\begin{lstlisting}
make all
\end{lstlisting}

This runs tests, compile checks, Week 8--12 demos, and the final demo. Docker is
also supported:

\begin{lstlisting}
make docker-build
make docker-run
\end{lstlisting}

The latest local verification passed with 42 tests. The Week 12 demo reported
10 registered features, 3 valid events, 2 intentionally invalid events, 6 batch
features written, and 1 lineage record. The final demo writes prediction logs,
lineage logs, and quality reports under \texttt{demo\_data/}. Generated files
are ignored by git and are meant for inspection, not as source code.

For the recommender experiment, the data source is explicit in the output. If
\path{data/raw/retailrocket/events.csv} exists, the run uses the RetailRocket
Kaggle events. Otherwise the output says that the source is the built-in sample.
In the current repository state the built-in sample is the reproducible default,
while the Kaggle path documents how to repeat the same pipeline with a larger
external ecommerce dataset.

\begin{center}
\footnotesize
\setlength{\tabcolsep}{4pt}
\renewcommand{\arraystretch}{0.95}
\begin{tabular}{p{0.25\linewidth} p{0.28\linewidth} p{0.36\linewidth}}
\toprule
\textbf{Check} & \textbf{Observed result} & \textbf{Meaning} \\
\midrule
Test suite & 42 tests passed & Core behavior is covered by automated checks. \\
Week 8 evaluation & JSON-lines metrics written & Runs are inspectable and reproducible. \\
Week 11 recommender & Model artifact and prediction log & The AI layer can consume ordered behavior. \\
Week 12 inspection & 10 registered features and 1 lineage record & Feature meanings and prediction inputs are traceable. \\
Docker run & Tests and demos pass in container & The project is easier to run outside my local machine. \\
\bottomrule
\end{tabular}
\vspace{0.2cm}
\textit{Representative local verification results.}
\end{center}

These results should be interpreted as functional validation, not as a claim of
production performance. The project checks that the layers work together, that
records are written where expected, and that generated artifacts can be
inspected. The absolute timing values are laptop-dependent and are less
important than the fact that the system exposes them.

\section{Comparison With Production Systems}\label{sec:comparison}

\begin{table}[!htbp]
\centering
\footnotesize
\setlength{\tabcolsep}{4pt}
\renewcommand{\arraystretch}{0.95}
\begin{tabular}{p{0.30\linewidth} p{0.25\linewidth} p{0.34\linewidth}}
\toprule
\textbf{MLStore-Lite layer} & \textbf{Reference} & \textbf{Shared idea} \\
\midrule
Storage engine & RocksDB / LevelDB & WAL, memtable, SSTables, compaction \\
Sharding and replication & Cassandra & partitioned key space, replicated data \\
Event log & Kafka & append-only events, producers, consumers, offsets \\
Batch processing & Spark / MapReduce & map, shuffle, reduce over finite data \\
Stream processing & Flink / Kafka Streams & continuous event processing and windows \\
Integrated workflow & Databricks / platforms & end-to-end data and ML infrastructure \\
Inference and recommender & feature stores / recommenders & feature retrieval, prediction logging, event sequences \\
Metadata and quality & ML observability tools & feature definitions, checks, lineage \\
\bottomrule
\end{tabular}
\caption{Conceptual comparison with production systems.}
\end{table}

The main difference is scale and operational complexity. Production systems run
across real machines, handle concurrent users, recover from many failure modes,
and optimize for performance. MLStore-Lite runs locally and prioritizes
readability.

This comparison was useful because it gave each local component a reference
point. For example, the storage engine made LSM-tree terminology concrete, the
stream layer made offsets and event logs concrete, and the recommender showed
why feature stores exist in ML workflows. I did not download or run these
production systems. They were used as architectural comparisons, while the
implemented system stayed small enough to understand directly.

\section{Design Tradeoffs}\label{sec:tradeoffs}

Several choices were made to keep the project educational. First, all nodes are
local Python objects rather than real network services. This removes network
setup, but keeps the concepts of nodes, replicas, leaders, and shards visible.
Second, the project uses simple JSON and local files rather than external
databases. This makes recovery logs, experiment logs, and prediction logs easy
to open in a text editor.

Third, the AI models are intentionally lightweight. The purchase-intent scorer
is deterministic, and the sequential recommender is a small attention model
rather than a full PyTorch training stack. This keeps the focus on the
infrastructure question: how do events become features, how are features served,
and how can predictions be traced? A larger model could improve ML realism, but
it would also make the repository harder to inspect.

Fourth, the project favors explicit scripts over hidden automation. The
Makefile provides one-command execution, but the underlying scripts still map to
the weekly build-up. This makes it possible to run the whole system or study
one layer at a time.

\section{Repository Map}\label{sec:repomap}

The repository is organized so that each conceptual layer has a small code
area. The storage code lives under \path{src/mlstore_lite/storage/}. This is
where the WAL, memtable, SSTable-like files, compaction, and recovery logic are
implemented. The replication code lives under
\path{src/mlstore_lite/replication/}, where nodes are grouped into leader-
follower clusters. The sharding code lives under
\path{src/mlstore_lite/sharding/}, where the consistent hash ring and sharded
cluster wrapper route keys to the correct replica group.

The data processing layers are split into batch and stream folders. The batch
engine under \path{src/mlstore_lite/batch/} computes features from finite event
lists. The stream code under \path{src/mlstore_lite/stream/} stores append-only
events, tracks consumer offsets, and computes windowed feature updates. The
integration layer connects these parts so that batch and stream outputs go into
the same sharded replicated store.

The AI and inspection layers are also separated. The feature-serving and model
logic live under \path{src/mlstore_lite/ai/}. Training helpers for the
RetailRocket-style event data live under \path{src/mlstore_lite/training/}.
The metadata, lineage, and quality checks live under \path{src/mlstore_lite/catalog/}
and related modules. Finally, \path{src/mlstore_lite/experiments/} contains the
runnable demos for evaluation, inference, recommendation, and the final
end-to-end run.

The tests under \path{tests/} are important because they keep the project from
being only a narrative report. They check storage behavior, replication,
sharding, batch processing, stream processing, observability, AI inference, and
the later inspection layer. The documentation under \path{docs/weekly-notes/}
acts as a learning diary, while this final report gives the compact submission
version of the same story.

\section{Limitations and Future Work}\label{sec:limitations}

The most important limitation is that the project is local. Nodes are Python
objects and directories, not independent networked processes. The project does
not include automatic consensus, real leader election, background repair,
transaction isolation, a model registry, HTTP serving, or production monitoring.
The AI layer is intentionally small, and the recommender is not trained as a
large neural network.

Future work could revisit the repository using the newer DDIA edition, add
networked nodes, implement a small consensus experiment, or build mini-projects
inspired by \textit{System Design Interview}. On the AI side, the tiny attention
model could be replaced by a PyTorch Transformer, a two-tower retrieval model,
or a Mamba-style state-space model for longer histories.

\section{Conclusion}\label{sec:conclusion}

MLStore-Lite implements a compact data infrastructure prototype for ML-style
features. It starts from a single-node storage engine and gradually adds
replication, sharding, batch processing, stream processing, observability,
feature serving, sequence-based recommendation, and inspection tools.

The final architecture can be summarized as:

\begin{lstlisting}
events -> quality checks -> features -> sharded store
       -> inference/recommender -> lineage -> prediction logs
\end{lstlisting}

The value of the project is not raw performance. The value is that the internals
are small enough to inspect, while still connecting to the architecture of real
data-intensive and ML infrastructure systems.

The main learning outcome is that abstract systems ideas become much clearer
when implemented end to end. Storage, replication, sharding, streams, and model
serving are often discussed as separate topics, but in an ML system they depend
on each other. The project helped me see that ML infrastructure is not only
about choosing a model. It is also about how data is written, copied, routed,
computed, served, checked, and explained after a prediction is made.
